% Part: turing-machines
% Chapter: machines-computations
% Section: turing-machines

\documentclass[../../../include/open-logic-section]{subfiles}

\begin{document}

\olfileid[pt]{tur}{mac}{tur}
\olsection{Máquinas de Turing}

\begin{explain}
A definição formal do que constitui uma máquina de Turing parece
abstrata, mas é na verdade simples: apenas reúne em uma
estrutura matemática toda a informação necessária para especificar o
funcionamento de uma máquina de Turing. Isso inclui (1) em quais estados a
máquina pode estar, (2) quais símbolos são permitidos na fita, (3)
em qual estado a máquina deve começar e (4) qual é o conjunto de instruções
da máquina.
\end{explain}

\begin{defn}[Máquina de Turing]
Uma \emph{máquina de Turing} $M$ é uma tupla $\langle Q, \Sigma, q_0,
\delta\rangle$ consistindo de
\begin{enumerate}
\item um conjunto finito de \emph{estados}~$Q$,
\item um \emph{alfabeto} finito $\Sigma$ que inclui $\TMendtape$ e
  $\TMblank$,
\item um \emph{estado inicial}~$q_0 \in Q$,
\item um \emph{conjunto de instruções} finito~$\delta\colon Q \times \Sigma
  \pto Q \times \Sigma \times \{\TMleft, \TMright, \TMstay\}$.
\end{enumerate}
A função parcial~$\delta$ também é chamada de \emph{função de transição}
de~$M$.
\end{defn}

\begin{explain}
Supomos que a fita é infinita em apenas uma direção. Por essa
razão é útil designar um símbolo especial~$\TMendtape$ como
marcador da extremidade esquerda da fita. Isso facilita que programas de
máquinas de Turing saibam quando estão ``em perigo'' de sair da
fita. Poderíamos supor que esse símbolo nunca é sobrescrito, ou seja,
que $\delta(q,\TMendtape) = \tuple{q', \TMendtape, x}$ se
$\delta(q,\TMendtape)$ estiver definido. Alguns livros-texto fazem isso; nós não.
Basta ter cuidado ao construir sua máquina de Turing para que
ela nunca sobrescreva~$\TMendtape$. Além disso, há casos em que
permitir tal sobrescrita oferece alguma flexibilidade conveniente.
\end{explain}

\begin{ex}
\emph{Máquina dos pares:} A máquina dos pares é formalmente
a quádrupla $\tuple{Q, \Sigma, q_0, \delta}$ onde
\begin{align*}
Q & = \{ q_0, q_1 \} \\
\Sigma & = \{ \TMendtape, \TMblank, \TMstroke \}, \\
\delta(q_0, \TMstroke) & = \tuple{q_1, \TMstroke, \TMright},\\
\delta(q_1, \TMstroke) & = \tuple{q_0, \TMstroke, \TMright},\\
\delta(q_1, \TMblank)  & = \tuple{q_1, \TMblank, \TMright}.
\end{align*}
\end{ex}

\end{document}
